\chapter{Sensitivity when the identifying assumptions are uncertain}
\label{r3:sensitivity}
The preceding estimators are conditional conclusions: given the
stated graph, moment restrictions, support and observation process,
they estimate a specified causal functional. A bank rarely knows
those premises with certainty. Sensitivity analysis should therefore
describe an explicit larger model class and propagate its ambiguity
to the financial target. Replacing one flexible regression by another
mostly probes estimation choices; it does not probe the omitted
mechanisms those regressions share.

\section{Hidden confounding as a difference between two functionals}
\label{r3:ovb-sensitivity}
Let the full covariates be $W=(A,X,U)$ and the observed ones
$W_s=(A,X)$. Write $g(W)=\E[Y\mid W]$ and
$g_s(W_s)=\E[Y\mid W_s]$. Suppose the desired parameter is a
continuous linear functional of the full regression, denoted
$\theta=\mathcal L(g)$, with square-integrable representer $\alpha$:
\[
 \mathcal L(f)=\E[\alpha(W)f(W)].
\]
The corresponding observed-data functional uses $g_s$ and
$\alpha_s(W_s)=\E[\alpha(W)\mid W_s]$; denote it by $\theta_s$.
For a binary-treatment ATE, full exchangeability given $(X,U)$ and
sufficient overlap supply
\[
 \alpha=\frac{A}{e(X,U)}-\frac{1-A}{1-e(X,U)},\qquad
 \alpha_s=\frac{A}{e_s(X)}-\frac{1-A}{1-e_s(X)}.
\]
Bayes' rule verifies the projection identity. The second expression
is computable from recorded covariates, but its associated $\theta_s$
need not be causal when $U$ is omitted. Representer existence itself
requires the continuity and moment conditions in
\citet[Assumptions 1--2 and Lemma 1]{ovb2022}; for instance weak
overlap with infinite inverse-propensity moments can violate them.

Orthogonal projection gives the signed identity
\begin{equation}
 \theta-\theta_s
    =\E[(g-g_s)(\alpha-\alpha_s)],\qquad
 |\theta-\theta_s|^2
    \le\E[(g-g_s)^2]\E[(\alpha-\alpha_s)^2].
 \label{r3:ovb-bound}
\end{equation}
To prove it, expand $\E[g\alpha]-\E[g_s\alpha_s]$ after
writing each full function as its projection plus residual.
The two cross terms vanish by conditional mean zero; Cauchy--Schwarz
bounds the remaining product. This is the expansion in
\citet[Appendix B.4]{ovb2022}.\footnote{The first signed equation
in Theorem 2 of the archived 2022 working paper reverses the
left-hand difference relative to its Appendix B.4 expansion.
Equation~\eqref{r3:ovb-bound} fixes the sign by the displayed
projection proof. The source's squared bound is unchanged.}

The bound has a useful calibration in observed financial units:
\begin{align}
 |\theta-\theta_s|&\le S C_Y C_A,\nonumber\\
 S^2&=\E[(Y-g_s)^2]\E[\alpha_s^2],\nonumber\\
 C_Y^2&=\frac{\E[(g-g_s)^2]}{\E[(Y-g_s)^2]},\qquad
 C_A^2=\frac{\E[\alpha^2]-\E[\alpha_s^2]}{\E[\alpha_s^2]}.
 \label{r3:ovb-calibration}
\end{align}
Assume the displayed denominators are positive. The observable scale
$S$ multiplies two unknown sensitivity parameters. $C_Y^2\in[0,1]$
measures how much residual outcome variation omitted variables explain;
$C_A^2\ge0$ measures their additional representer variation and
need not be at most one. These are the nonlinear functional quantities
in \citet[Corollary 2]{ovb2022}, not automatically the partial
$R^2$ values of a chosen linear regression. The covariance's
correlation can tighten the bound if a restriction on its sign or
magnitude is defensible.

For a finite sign and magnitude check, use the hidden-confounding
model from Section~\ref{r3:nonidentification}: $U$ is fair,
$P(A=1\mid U)=.2+.6U$, and $P(Y=1\mid A,U)=.2+.6U$.
The true ATE is zero and the observed short contrast is $.36$.
Direct enumeration of the projection product gives $-.36$, in
agreement with $\theta-\theta_s$. The squared upper bound is
$.1296$, attained in this construction. The example checks the
signed identity and a sharp special case; it does not prove that
every imposed sensitivity pair is compatible with every observed law.
The source's sharpness statement is conditional on admissible
function choices without additional observed-law constraints.

Measured-covariate deletion can supply a benchmark: compare the
estimated regression and representer changes when removing an
observed variable of known operational importance. Assuming an
unmeasured variable is no stronger than that benchmark is still a
substantive judgment. Better prediction of recorded outcomes does
not demonstrate small $C_Y$ or $C_A$. Statistical intervals for
the learnable scale and short effect also need to account for their
joint estimation; merely appending the point estimate of a bias
bound to an ordinary interval is not generally a calibrated procedure.

\section{Trend restrictions produce an identified set}
\label{r3:trend-sensitivity}
Let $\beta$ be a vector of population event-study contrasts, all
relative to a reference date normalized to zero. With no anticipation,
write
\[
 \beta_{\rm pre}=\delta_{\rm pre},\qquad
 \beta_{\rm post}=\tau_{\rm post}+\delta_{\rm post}.
\]
Here $\delta$ describes counterfactual differences between groups in
the absence of treatment and $\tau$ is the causal path effect. The
target is $\theta=l'\tau_{\rm post}$ for declared weights $l$.
Rather than impose exact parallel trends, choose a restriction set
$\Delta$ for $\delta$. Algebra alone gives
\begin{equation}
 \Theta(\beta,\Delta)=
 \{l'\beta_{\rm post}-l'\delta_{\rm post}:
                \delta\in\Delta,\ \delta_{\rm pre}=\beta_{\rm pre}\}.
 \label{r3:trend-set}
\end{equation}
For closed convex restrictions with finite attained extrema,
the endpoints are obtained by maximizing and minimizing
$l'\delta_{\rm post}$ under those constraints. Unbounded sets
produce infinite endpoints; an empty feasible set indicates a
restriction inconsistent with that population vector. This is the
construction in \citet[Section 2.3 and Lemma 2.1]{rambachan2023}.

Two scientifically different choices illustrate why a sensitivity
parameter needs a definition. A relative-magnitude restriction uses
\[
 |\delta_{t+1}-\delta_t|
 \le\overline M\max_{s<0}|\delta_{s+1}-\delta_s|\quad(t\ge0).
\]
A smoothness restriction instead uses
\[
 |\delta_{t+1}-2\delta_t+\delta_{t-1}|\le M.
\]
The first limits future deviations by past deviations. The second
limits changes in the differential trend's slope; $M=0$ continues
a linear differential trend rather than enforcing a zero one.
Both definitions and their interpretations are given in
\citet[Sections 2.4.1--2.4.4]{rambachan2023}.

For one pretreatment contrast $\beta_{-1}=-1$, baseline
$\delta_0=0$ and one post-treatment contrast $\beta_1=3$,
relative magnitude gives $\theta\in[3-\overline M,3+\overline M]$.
Smoothness gives $\delta_1\in[1-M,1+M]$ and hence
$\theta\in[2-M,2+M]$. These are population identification intervals
under different economic restrictions, not alternative standard
errors around the same point-identified effect. A branch with a
continuing adverse local demand trend may justify a different set
from one exposed to a new abrupt competitor entry.

\section{Sampling uncertainty and assumption uncertainty must be combined}
\label{r3:sensitivity-inference}
Replacing $\beta$ by $\widehat\beta$ in
Equation~\eqref{r3:trend-set} supplies estimated endpoints. It does
not by itself supply a confidence interval, because pretreatment
contrasts used in the restrictions are noisy and the endpoints can
be nonsmooth functions of them. A transparent conservative
construction starts with a valid joint confidence region
$\mathcal C_{1-\alpha}$ for the entire $\beta$ vector and returns
\[
 \bigcup_{b\in\mathcal C_{1-\alpha}}\Theta(b,\Delta).
\]
On the event $\beta\in\mathcal C_{1-\alpha}$, this set contains
the true $\theta$ whenever the stated restriction is correct.
That inclusion proves its coverage. The region must respect the
cluster/time dependence of the design, and an asymptotic normal
region inherits its approximation limits. The conditional,
hybrid and fixed-length methods in
\citet[Sections 3--4]{rambachan2023} offer more powerful procedures
under their own regularity and restriction classes; they are not
equivalent to endpoint substitution.

Passing a pretreatment trend test cannot establish the extrapolation
restriction. Low power leaves large violations plausible, and using
the same test to select which studies or specifications to report
changes the sampling problem. A useful banking analysis reports
which financial conclusion changes as $M$ or $\overline M$ grows,
why those departures are operationally plausible, and which new
measurement or experiment would narrow them.

The same pattern applies to invalid-IV ranges, alternative discovery
graphs, uncertain entity linkage and model-dependent POMDP values.
Each enlarges a different part of the causal model class. Their
combination must be defined jointly; independently varying several
knobs and taking the largest separate effect need not bound their
combined effect. A policy recommendation can be robust across a
declared ambiguity set, or conditional on one defended model. Neither
label should hide how that model or set was chosen.
For the linkage problem, Section~\ref{s61:linkage-sets} executes this
confidence-region inclusion argument through a finite interval-union
algorithm. Its full-line results also show why coverage and decision
informativeness must be reported separately.
